\documentclass[12pt,a4paper]{article}
\usepackage[utf8]{inputenc}
\usepackage[T1]{fontenc}
\usepackage[polish]{babel}
\usepackage{amsmath,amssymb}
\usepackage{graphicx}
\usepackage{booktabs}
\usepackage{hyperref}
\usepackage{geometry}
\usepackage{caption}
\usepackage{subcaption}
\usepackage{float}
\usepackage{enumitem}
\usepackage{csquotes}
\usepackage[backend=biber,style=numeric,sorting=nyt]{biblatex}
\addbibresource{references.bib}

\geometry{margin=2.5cm}

\title{Zastosowanie Monte Carlo Tree Search do gry w warcaby \\[0.5em]
\large Raport końcowy — Metody Sztucznej Inteligencji 2}
\author{
  % TODO: Wpisać imiona i nazwiska autorów
  Autor 1 \and Autor 2
}
\date{\today}

\begin{document}
\maketitle

% ===========================================================================
\begin{abstract}
W niniejszym raporcie przedstawiono badanie skuteczności algorytmu Monte Carlo Tree Search (MCTS) z~formułą Upper Confidence Bounds Applied to Trees (UCT) w~grze w~warcaby na planszy 8\texttimes8 z~damkami latającymi. Porównano cztery warianty graczy komputerowych: czysty UCT, UCT z~płytkim przeszukiwaniem minimax (hybryda wg Baier i~Winands, 2014), UCT z~wczesnym obcinaniem symulacji i~ewaluacją heurystyczną (Lorentz, 2016; Lanctot et~al., 2014) oraz gracz heurystyczny oparty o~alpha-beta z~ewaluacją inspirowaną programem Chinook. Zbadano wpływ stałej eksploracji~$C$, budżetu iteracji oraz limitu ruchów na skuteczność poszczególnych podejść. Eksperymenty przeprowadzono na co~najmniej 50~partiach dla każdej konfiguracji z~wielokrotnymi ziarnami generatora losowego w~celu zapewnienia powtarzalności i~oceny rozrzutu wyników. Wyniki wskazują, że % TODO: uzupełnić po eksperymentach
zwiększenie budżetu iteracji istotnie poprawia skuteczność UCT, a~warianty hybrydowe oferują mniejszą wariancję wyników kosztem dłuższego czasu obliczeń. Skrócenie limitu ruchów zmienia ranking algorytmów, faworyzując bardziej agresywne strategie materiałowe.
\end{abstract}

% ===========================================================================
\section{Wprowadzenie}
\label{sec:intro}

Monte Carlo Tree Search (MCTS) to rodzina algorytmów przeszukiwania drzewa gry, która od czasu przełomowych wyników w~Go \cite{silver2016mastering} znajduje zastosowanie w~wielu grach planszowych i~problemach decyzyjnych. Kluczową zaletą MCTS jest brak konieczności definiowania funkcji oceny pozycji — algorytm uczy się wartości stanów z~losowych symulacji (playoutów). Wariant UCT \cite{kocsis2006bandit} łączy przeszukiwanie drzewa z~formułą UCB1, zapewniając asymptotyczną zbieżność do optymalnej strategii.

Celem niniejszej pracy jest zbadanie skuteczności UCT i~jego wariantów w~grze w~warcaby (draughts) na planszy 8\texttimes8 z~damkami latającymi. Warcaby są grą o~pełnej informacji z~przestrzenią stanów rzędu $5 \times 10^{20}$ \cite{schaeffer2007checkers}, co czyni je odpowiednim testbedem — gra jest wystarczająco złożona, by wymagać zaawansowanych algorytmów, a~jednocześnie istnieją silne programy heurystyczne (Chinook) jako punkt odniesienia.

\subsection{Cel badania}
Porównanie czterech graczy komputerowych pod kątem:
\begin{itemize}[nosep]
  \item skuteczności (win rate) przy różnych budżetach obliczeniowych,
  \item stabilności wyników (wariancja między ziarnami RNG),
  \item czasu obliczeń na ruch,
  \item wrażliwości na zmianę parametrów gry (limit ruchów).
\end{itemize}

% ===========================================================================
\section{Opis problemu}
\label{sec:problem}

\subsection{Zasady gry}
Zastosowano wariant warcabow na planszy 8\texttimes8 z~następującymi regułami:
\begin{itemize}[nosep]
  \item Plansza 8\texttimes8, ruch wyłącznie po ciemnych polach (32 pola gry).
  \item 12 pionków na stronę; białe zaczynają (na rzędach 5--7), czarne na rzędach 0--2.
  \item Białe ruszają się jako pierwsze.
  \item Pionek porusza się o~jedno pole po przekątnej do przodu.
  \item Bicie jest \textbf{obowiązkowe} — jeśli istnieje ruch bijący, gracz musi bić.
  \item Pionek bije do przodu \textbf{i} do tyłu (skok nad sąsiednim wrogiem na puste pole za nim).
  \item Wielobicia są obowiązkowe i~kontynuowane tym samym pionkiem aż do wyczerpania możliwości.
  \item Zbite figury usuwane są dopiero po zakończeniu całej sekwencji bicia.
  \item Damki są \textbf{latające}: poruszają się i~biją po całej przekątnej (mogą przeskoczyć wroga na dowolnej odległości i~wylądować na dowolnym pustym polu za nim).
  \item Promocja piona do damki następuje po zakończeniu całego ruchu na ostatnim rzędzie.
  \item Limit ruchów: 150 półruchów (plies); po przekroczeniu — remis.
  \item Zasada maksymalnego bicia jest \textbf{wyłączona} (dowolna pełna sekwencja bicia jest legalna).
\end{itemize}

\subsection{Modyfikacje względem klasycznych zasad}
\begin{itemize}[nosep]
  \item Wprowadzono limit 150 ruchów w~celu uniknięcia nieskończonych partii (standardowe zasady nie przewidują takiego limitu).
  \item Zasada najdłuższego bicia została wyłączona — upraszcza implementację i~nie zmienia istotnie charakteru gry.
  \item Promocja w~trakcie wielobicia następuje dopiero na końcu sekwencji (pion kontynuuje jako pion nawet jeśli przebiega przez rząd promocji).
\end{itemize}

% ===========================================================================
\section{Przegląd literatury}
\label{sec:literature}

\subsection{MCTS i UCT}
Algorytm MCTS w~wariancie UCT został zaproponowany przez Kocsisa i~Szepesváriego \cite{kocsis2006bandit}. Formuła selekcji:
\begin{equation}
  a^* = \arg\max_a \left[ \frac{Q(s,a)}{N(s,a)} + C \cdot \sqrt{\frac{\ln N(s)}{N(s,a)}} \right]
  \label{eq:ucb1}
\end{equation}
gdzie $Q(s,a)$ to skumulowana nagroda, $N(s,a)$ liczba odwiedzin, a~$C$ stała eksploracji. Świechowski et~al. \cite{swiechowski2023mcts} przedstawiają obszerny przegląd modyfikacji MCTS.

\subsection{UCT + Minimax (hybryda)}
Baier i~Winands \cite{baier2014mcts} zaproponowali zastąpienie losowego playoutu płytkim przeszukiwaniem minimax. Zamiast symulacji do końca gry, z~liścia drzewa MCTS uruchamiany jest minimax o~głębokości~$d$, który zwraca dokładniejszą ocenę pozycji. Zmniejsza to wariancję wyników kosztem większego czasu obliczeniowego na iterację.

\subsection{UCT + Early Termination}
Lorentz \cite{lorentz2016evaluation} zaproponował obcinanie symulacji po ustalonej liczbie ruchów i~ocenę pozycji heurystyką. Lanctot et~al. \cite{lanctot2014mcts} wprowadzili blending:
\begin{equation}
  \hat{Q} = (1 - \alpha) \cdot \frac{r}{n} + \alpha \cdot v
  \label{eq:blending}
\end{equation}
gdzie $v$ to ewaluacja heurystyczna węzła, a~$\alpha \in [0,1]$ kontroluje siłę blendingu.

\subsection{Heurystyki w warcabach}
Program Chinook \cite{schaeffer1997one} stosował ewaluację łączącą materiał (pionek=1, damka=3), kontrolę centrum, mobilność i~bliskość promocji. Podejście to stanowi nasz benchmark heurystyczny.

% ===========================================================================
\section{Metodologia}
\label{sec:methodology}

\subsection{Zaimplementowani gracze}

\subsubsection{UCT bazowy}
Standardowy cztero-fazowy MCTS: selekcja (UCB1, formuła~\ref{eq:ucb1}), ekspansja (jeden losowy nieodwiedzony ruch), symulacja (losowy playout do stanu terminalnego), propagacja wsteczna. Parametry: liczba iteracji, stała~$C$, ziarno RNG.

\subsubsection{UCT + Minimax}
Jak UCT bazowy, ale faza symulacji zastąpiona przeszukiwaniem minimax z~obcinaniem alpha-beta o~konfigurowalnej głębokości~$d$. Funkcja oceny liścia: heurystyka materiałowa znormalizowana do $[-1, 1]$:
\begin{equation}
  h = \frac{M_{\text{own}} - M_{\text{opp}}}{M_{\text{own}} + M_{\text{opp}}}
  \label{eq:material}
\end{equation}
gdzie $M = \text{(liczba pionków)} \cdot 1 + \text{(liczba damek)} \cdot 3$.

\subsubsection{UCT + Early Termination}
Symulacja jest obcinana po $k$ ruchach; pozycja oceniana heurystyką materiałową (formuła~\ref{eq:material}). Opcjonalnie stosowany blending (formuła~\ref{eq:blending}) z~parametrem~$\alpha$.

\subsubsection{Alpha-beta heurystyczny}
Iterative-deepening alpha-beta (negamax) z~ewaluacją inspirowaną Chinookiem:
\begin{itemize}[nosep]
  \item materiał (pionek=1, damka=3),
  \item bonus za kontrolę centrum (pola 3--4 w~obu osiach),
  \item bliskość promocji (im bliżej rzedu promocji, tym wyżej),
  \item bonus za obronę ostatniego rzędu.
\end{itemize}
Maksymalna głębokość przeszukiwania: parametr \texttt{max\_depth} (domyślnie 6).

\subsection{Powtarzalność}
Każdy gracz oparty na losowości przyjmuje jawne ziarno generatora (parametr \texttt{seed}). Dla tego samego ziarna wybór ruchu jest deterministyczny. Każdy eksperyment loguje użyte ziarno do pliku CSV.

\subsection{Plan eksperymentów}
Eksperymenty zorganizowane w~macierz 18 konfiguracji (Tabela~\ref{tab:experiments}):
\begin{itemize}[nosep]
  \item \textbf{H1}: UCT vs AlphaBeta — sweep po $C \in \{0.1, 0.3, 0.7, 1.0, 1.4, 2.0, 3.0\}$ przy 100 iteracjach; sweep po budżecie $\{100, 500, 1000, 5000, 10000\}$ przy $C=1.4$.
  \item \textbf{H2}: UCT+Minimax vs UCT — 1000 i 5000 iteracji, $d=2$.
  \item \textbf{H3}: UCT+EarlyTerm vs UCT — 1000 i 5000 iteracji, $k=20$, $\alpha=0.5$.
  \item \textbf{H4}: Powtórzenie wybranych porównań z~limitem 100 ruchów (vs domyślne 150).
\end{itemize}

Dla każdej konfiguracji: minimum 50 partii, minimum 20 ziaren, zamiana kolorów (co druga partia).

\begin{table}[H]
\centering
\caption{Macierz eksperymentów}
\label{tab:experiments}
\small
\begin{tabular}{llllrrrl}
\toprule
ID & Hipoteza & Gracz A & Gracz B & Iteracje & $C$ & Limit & Uwagi \\
\midrule
E01--E07 & H1 & UCT & AlphaBeta & 100 & 0.1--3.0 & 150 & C-sweep \\
E08--E11 & H1 & UCT & AlphaBeta & 500--10k & 1.4 & 150 & Budget-sweep \\
E12--E13 & H2 & UCT+Mini & UCT & 1k/5k & 1.4 & 150 & $d=2$ \\
E14--E15 & H3 & UCT+ET & UCT & 1k/5k & 1.4 & 150 & $k=20, \alpha=0.5$ \\
E16--E18 & H4 & Różni & AlphaBeta & 1000 & 1.4 & 100 & Skrócony limit \\
\bottomrule
\end{tabular}
\end{table}

% ===========================================================================
\section{Hipotezy badawcze}
\label{sec:hypotheses}

\begin{description}
  \item[H1] Czysty UCT, po strojeniu parametru $C$, osiąga wyższy win rate od gracza heurystycznego dla niskiego i~średniego budżetu iteracji.
  \item[H2] UCT + minimax rollout (płytki) obniża wariancję wyników między ziarnami względem czystego UCT, oferując bardziej stabilną jakość decyzji.
  \item[H3] UCT + early termination z~ewaluacją heurystyczną skraca średni czas na ruch przy porównywalnym win rate do czystego UCT.
  \item[H4] Zmiana limitu ruchów ze 150 na 100 zmienia ranking algorytmów i~zwiększa odsetek zwycięstw graczy agresywnie materialowych.
\end{description}

% ===========================================================================
\section{Wyniki eksperymentów}
\label{sec:results}

% TODO: Uzupełnić po przeprowadzeniu pełnych eksperymentów.
% Poniżej przygotowane szablony na wykresy i tabele.

\subsection{H1: Wpływ stałej eksploracji $C$}

\begin{figure}[H]
\centering
% \includegraphics[width=0.8\textwidth]{results/plots/h1_c_sweep.png}
\caption{Win rate UCT vs AlphaBeta w~funkcji stałej eksploracji $C$ (100 iteracji, 50 partii na punkt). Słupki błędów: 95\% CI.}
\label{fig:c_sweep}
\end{figure}

% TODO: Opis wyników — przy jakim C UCT osiąga najlepsze rezultaty?

\subsection{H1: Wpływ budżetu iteracji}

\begin{figure}[H]
\centering
% \includegraphics[width=0.8\textwidth]{results/plots/h1_budget_sweep.png}
\caption{Win rate UCT vs AlphaBeta w~funkcji budżetu iteracji ($C=1.4$, 50 partii na punkt).}
\label{fig:budget_sweep}
\end{figure}

% TODO: Przy jakim budżecie UCT zaczyna dominować nad alpha-beta?

\subsection{H2 i H3: Porównanie wariantów UCT}

\begin{figure}[H]
\centering
% \includegraphics[width=0.8\textwidth]{results/plots/h2_h3_variants.png}
\caption{Win rate wariantów UCT (minimax, early termination) vs czysty UCT.}
\label{fig:variants}
\end{figure}

\begin{table}[H]
\centering
\caption{Porównanie wariancji win rate między ziarnami}
\label{tab:variance}
\begin{tabular}{lrrr}
\toprule
Gracz & Win rate & Std & Avg time/move [s] \\
\midrule
UCT (1000 iter) & -- & -- & -- \\
UCT+Minimax (1000 iter, $d=2$) & -- & -- & -- \\
UCT+EarlyTerm (1000 iter, $k=20$) & -- & -- & -- \\
\bottomrule
\end{tabular}
\end{table}

% TODO: Uzupełnić tabelę danymi z eksperymentów E12-E15.

\subsection{H4: Wpływ limitu ruchów}

\begin{figure}[H]
\centering
% \includegraphics[width=0.8\textwidth]{results/plots/h4_move_limit.png}
\caption{Porównanie win rate vs AlphaBeta przy limicie 150 i~100 ruchów.}
\label{fig:move_limit}
\end{figure}

% TODO: Czy skrócenie limitu zmienia ranking? Kto zyskuje, kto traci?

\subsection{Czas obliczeń}

\begin{figure}[H]
\centering
% \includegraphics[width=0.9\textwidth]{results/plots/time_per_move.png}
\caption{Średni czas na ruch dla poszczególnych eksperymentów.}
\label{fig:time}
\end{figure}

% ===========================================================================
\section{Testy człowiek--komputer}
\label{sec:human_tests}

Interfejs graficzny zrealizowano w~bibliotece Pygame. Przeprowadzono testy z~udziałem co najmniej 5 osób. Każdy uczestnik rozegrał partię z~wybranym botem i~wypełnił krótki formularz ewaluacyjny.

\begin{table}[H]
\centering
\caption{Wyniki testów człowiek--komputer}
\label{tab:human}
\begin{tabular}{clcccl}
\toprule
\# & Przeciwnik & Wynik & Czytelność GUI & Siła bota & Uwagi \\
\midrule
1 & -- & -- & -- & -- & -- \\
2 & -- & -- & -- & -- & -- \\
3 & -- & -- & -- & -- & -- \\
4 & -- & -- & -- & -- & -- \\
5 & -- & -- & -- & -- & -- \\
\bottomrule
\end{tabular}
\end{table}

% TODO: Uzupełnić po testach z użytkownikami.

% ===========================================================================
\section{Wnioski}
\label{sec:conclusions}

\begin{description}
  \item[H1:] % TODO: Czy UCT z dobranym C pokonuje alpha-beta? Przy jakim budżecie?
  \item[H2:] % TODO: Czy hybryda minimax daje mniejszą wariancję?
  \item[H3:] % TODO: Czy early termination skraca czas przy zachowaniu win rate?
  \item[H4:] % TODO: Czy zmiana limitu ruchów zmienia ranking?
\end{description}

\subsection{Ograniczenia badania}
\begin{itemize}[nosep]
  \item Głębokość minimax w~hybrydzie ograniczona do $d=2$ ze względu na czas obliczeń.
  \item Alpha-beta z~głębokością 6 może nie być optymalnym benchmarkiem dla najwyższych budżetów UCT.
  \item Brak optymalizacji silnika (np.\ bitboardy) ogranicza maksymalną liczbę iteracji w~rozsądnym czasie.
\end{itemize}

\subsection{Kierunki dalszych badań}
\begin{itemize}[nosep]
  \item RAVE/AMAF jako dodatkowe usprawnienie UCT.
  \item Progressive widening w~drzewie MCTS.
  \item Learned evaluation (np.\ prosta sieć neuronowa) zamiast ręcznej heurystyki.
\end{itemize}

% ===========================================================================
\section{Wykorzystanie dużych modeli językowych}
\label{sec:llm}

W~ramach projektu wykorzystano duże modele językowe (LLM) w~następującym zakresie:

\subsection{Kod źródłowy}
\begin{itemize}[nosep]
  \item Generowanie szkieletu architektury projektu (struktura katalogów, interfejsy).
  \item Implementacja silnika reguł gry (generator ruchów, wielobicia, damki latające).
  \item Implementacja graczy (UCT, UCT+Minimax, UCT+EarlyTerm, AlphaBeta).
  \item Implementacja GUI (Pygame) i harnessu eksperymentalnego.
  \item Generowanie skryptów analizy statystycznej i wizualizacji.
\end{itemize}
Każdy wygenerowany fragment był weryfikowany przez autorów poprzez testy jednostkowe i~smoke testy runtime.

\subsection{Raport}
\begin{itemize}[nosep]
  \item Generowanie szablonu LaTeX raportu końcowego.
  \item Wsparcie przy formułowaniu opisów metodologii.
\end{itemize}
Treść merytoryczna (interpretacja wyników, wnioski) jest autorska.

% ===========================================================================
\printbibliography

\end{document}
